\begin{foreignabstract}

The Food Acquisition Program is a Brazilian federal government project that purchases food directly from family farmers at market prices, without the need for bidding processes, and distributes it to people experiencing food insecurity, as well as to the social assistance network, public facilities, and the public education system.
In order to manage the program in the PDS Osvaldo de Oliveira settlement, a web system was developed to cover all processes carried out by the farmers and the responsible manager. Software engineering approaches, such as requirements gathering and agile methodology, were used in the development of the system.
The project follows the MVC architectural pattern, with the module responsible for data processing and storage separated from the user interaction module in different containers. The project deployment was carried out using GitLab Actions to generate the image and send it to the staging server.
Docker containers were used to isolate the modules, and Docker Compose was used to start the application. The entire project is open source and available on GitLab.

\vspace{\baselineskip}
\noindent\textbf{Keywords:} Social Technology, Free Software, Community Management, Software Engineering.

\end{foreignabstract}

